\scanpage{161}
\mspage{110}

est appelée une $\unsure{W|\mathcal W}$-cofibration, si pour tte flèche $f:\deleted{\unsure{A}}U'\longrightarrow U''$ dans $U\backslash\mathrm{Cat}$, la flèche $g:X'\mathbin{-}X''$ qui s'en déduit par le coch. de base \deleted{\unsure{par|sur}} $X$ (objet de $U\backslash\mathrm{Cat}$) est dans $W$ quand $f$ l'est.

\uline{NB} Cette notion \inserted{de $W$-cofibration} est stable par composition et par cochangement de base, et les iso. de $\mathrm{Cat}$ sont des $W$-cofibrations. Si $f$ est une $W$-cofibration, $f^{\mathrm o}$ l'est et inversement.

[ Le corollaire s'énonce en disant qu'une imm. ouverte ou fermée qui est \inserted{$W$-}dwyérienne est une $W$-cofibration. ]

d) Une flèche $i:U\mathbin{-}X$ dans $\mathrm{Cat}$ est appelée une \uline{$W$-cofibration triviale}, ou une $W$-\uline{équivalence couniverselle}, si pour tout $U'$ dans $U\backslash\mathrm{Cat}$, i.e. \deleted{\unsure{tout}} flèche $\unsure{p|f}:U\mathbin{-}U'$ dans $\mathrm{Cat}$, la flèche déduite de $i$

\[
i':U'\mathbin{-}X'=U'\amalg_{U}X'
\]

est dans $W$. (NB Il est immédiat que cela \unsure{entraîne|entraîne,} bien que $i$ est une $W$-cofibration.) [ Cette notion de $W$-cofibration triviale est elle aussi stable par composition et par cochangement de base, et les iso de $\mathrm{Cat}$ en sont. ]
% archive sheet number: 161
